% ============================================================
\chapter{SQL --- Agr\'egation et Fonctions de Fen\^etre}
\label{ch:agregation}
% ============================================================

Jusqu'ici, nos requ\^etes SQL retournaient des lignes individuelles. Mais souvent, on veut des r\'esum\'es~: le nombre total d'\'etudiants, la moyenne des notes par d\'epartement, le chiffre d'affaires par trimestre. C'est le r\^ole des \emph{fonctions d'agr\'egation} (\texttt{COUNT}, \texttt{SUM}, \texttt{AVG}, \texttt{MAX}, \texttt{MIN}) combin\'ees avec \texttt{GROUP BY}. Les \emph{fonctions de fen\^etre} (\texttt{OVER}), introduites dans SQL:2003, vont plus loin~: elles calculent des agr\'egats tout en conservant les lignes individuelles --- rang, cumul, moyenne mobile --- sans n\'ecessiter de sous-requ\^etes complexes.

\section{Fonctions d'agr\'egation}

\begin{definition}[Fonction d'agr\'egation]
Une \emph{fonction d'agr\'egation} op\`ere sur un ensemble de lignes et retourne
une valeur unique r\'esumant cet ensemble.
\end{definition}

\begin{formules}[title=\textbf{Fonctions d'agr\'egation standard}]
\begin{center}
\begin{tabular}{lll}
\toprule
\textbf{Fonction} & \textbf{Description} & \textbf{Ignore NULL ?} \\
\midrule
\texttt{COUNT(*)} & Nombre de lignes & Non \\
\texttt{COUNT(col)} & Nombre de valeurs non NULL & Oui \\
\texttt{COUNT(DISTINCT col)} & Nombre de valeurs distinctes & Oui \\
\texttt{SUM(col)} & Somme & Oui \\
\texttt{AVG(col)} & Moyenne & Oui \\
\texttt{MIN(col)} & Minimum & Oui \\
\texttt{MAX(col)} & Maximum & Oui \\
\bottomrule
\end{tabular}
\end{center}
\end{formules}

\begin{codeblock}[title=\textbf{Agr\'egations simples}]
\begin{minted}{sql}
SELECT COUNT(*) AS nb_inscriptions,
       COUNT(DISTINCT id_etudiant) AS nb_etudiants,
       AVG(note) AS moyenne,
       MIN(note) AS note_min,
       MAX(note) AS note_max,
       SUM(note) AS somme_notes
FROM Inscriptions;
\end{minted}
\end{codeblock}

\begin{codeoutput}
\begin{verbatim}
nb_inscriptions | nb_etudiants | moyenne | note_min | note_max | somme_notes
----------------+--------------+---------+----------+----------+------------
8               | 5            | 14.375  | 9.0      | 18.0     | 115.0
\end{verbatim}
\end{codeoutput}

\begin{attention}[title=\textbf{AVG et les valeurs NULL}]
\texttt{AVG} ignore les valeurs \texttt{NULL}. La moyenne de $\{10, \texttt{NULL}, 20\}$
est $15$, pas $10$. Si vous souhaitez traiter les NULL comme $0$, utilisez
\texttt{AVG(COALESCE(note, 0))}.
\end{attention}

\section{GROUP BY --- Regroupement}

\begin{definition}[GROUP BY]
La clause \texttt{GROUP BY} partitionne les lignes en groupes selon une ou
plusieurs colonnes. Les fonctions d'agr\'egation sont alors appliqu\'ees
\`a chaque groupe s\'epar\'ement.
\end{definition}

\begin{formules}[title=\textbf{Ordre d'ex\'ecution avec GROUP BY}]
\begin{enumerate}
    \item \texttt{FROM} / \texttt{JOIN}
    \item \texttt{WHERE} (filtre avant regroupement)
    \item \texttt{GROUP BY}
    \item \texttt{HAVING} (filtre apr\`es regroupement)
    \item \texttt{SELECT}
    \item \texttt{DISTINCT}
    \item \texttt{ORDER BY}
    \item \texttt{LIMIT}
\end{enumerate}
\end{formules}

\begin{codeblock}[title=\textbf{Moyenne par \'etudiant}]
\begin{minted}{sql}
SELECT e.nom, e.prenom,
       COUNT(*) AS nb_cours,
       ROUND(AVG(i.note), 2) AS moyenne,
       MIN(i.note) AS note_min,
       MAX(i.note) AS note_max
FROM Etudiants e
JOIN Inscriptions i ON e.id_etudiant = i.id_etudiant
GROUP BY e.id_etudiant, e.nom, e.prenom
ORDER BY moyenne DESC;
\end{minted}
\end{codeblock}

\begin{codeoutput}
\begin{verbatim}
nom     | prenom | nb_cours | moyenne | note_min | note_max
--------+--------+----------+---------+----------+---------
Roux    | Emma   | 1        | 18.0    | 18.0     | 18.0
Moreau  | Claire | 2        | 15.75   | 14.0     | 17.5
Bernard | Alice  | 2        | 14.25   | 13.0     | 15.5
Petit   | Bob    | 2        | 14.0    | 12.0     | 16.0
Leroy   | David  | 1        | 9.0     | 9.0      | 9.0
\end{verbatim}
\end{codeoutput}

\begin{antipattern}[title=\textbf{Colonne hors GROUP BY sans agr\'egation}]
Toute colonne dans le \texttt{SELECT} doit soit appara\^itre dans le
\texttt{GROUP BY}, soit \^etre dans une fonction d'agr\'egation~:
\begin{minted}{sql}
-- ERREUR : prenom n'est ni dans GROUP BY ni agrege
SELECT nom, prenom, AVG(note) FROM Inscriptions
JOIN Etudiants ON ... GROUP BY nom;

-- CORRECT :
SELECT nom, prenom, AVG(note) FROM Inscriptions
JOIN Etudiants ON ... GROUP BY nom, prenom;
\end{minted}
MySQL en mode non strict accepte cette erreur silencieusement (r\'esultat
ind\'etermin\'e). PostgreSQL et SQLite la rejettent correctement.
\end{antipattern}

\section{HAVING --- Filtrage apr\`es regroupement}

\begin{definition}[HAVING]
La clause \texttt{HAVING} filtre les groupes apr\`es agr\'egation. Elle est
\`a \texttt{GROUP BY} ce que \texttt{WHERE} est \`a \texttt{FROM}~:
\texttt{WHERE} filtre les lignes, \texttt{HAVING} filtre les groupes.
\end{definition}

\begin{codeblock}[title=\textbf{Filtrage avec HAVING}]
\begin{minted}{sql}
-- Etudiants ayant une moyenne superieure a 14
SELECT e.nom, AVG(i.note) AS moyenne
FROM Etudiants e
JOIN Inscriptions i ON e.id_etudiant = i.id_etudiant
GROUP BY e.id_etudiant, e.nom
HAVING AVG(i.note) > 14
ORDER BY moyenne DESC;

-- Cours avec au moins 2 inscrits
SELECT c.intitule, COUNT(*) AS nb_inscrits
FROM Cours c
JOIN Inscriptions i ON c.id_cours = i.id_cours
GROUP BY c.id_cours, c.intitule
HAVING COUNT(*) >= 2;
\end{minted}
\end{codeblock}

\begin{codeoutput}
\begin{verbatim}
-- Moyenne > 14 :
nom     | moyenne
--------+--------
Roux    | 18.0
Moreau  | 15.75
Bernard | 14.25

-- Cours avec >= 2 inscrits :
intitule         | nb_inscrits
-----------------+------------
Bases de donnees | 3
Algorithmes      | 2
Statistiques     | 2
\end{verbatim}
\end{codeoutput}

\section{Groupements avanc\'es}

\subsection{GROUPING SETS, ROLLUP, CUBE}

\begin{codeblock}[title=\textbf{ROLLUP --- sous-totaux hi\'erarchiques (PostgreSQL)}]
\begin{minted}{sql}
SELECT e.filiere,
       c.intitule,
       AVG(i.note) AS moyenne,
       COUNT(*) AS nb
FROM Etudiants e
JOIN Inscriptions i ON e.id_etudiant = i.id_etudiant
JOIN Cours c ON i.id_cours = c.id_cours
GROUP BY ROLLUP (e.filiere, c.intitule)
ORDER BY e.filiere, c.intitule;
\end{minted}
\end{codeblock}

\texttt{ROLLUP} g\'en\`ere les sous-totaux hi\'erarchiques~: par (fili\`ere, cours),
par fili\`ere, et le total g\'en\'eral.

\texttt{CUBE} g\'en\`ere toutes les combinaisons possibles de sous-totaux.

\section{Fonctions de fen\^etre (Window Functions)}

\begin{definition}[Fonction de fen\^etre]
Une \emph{fonction de fen\^etre} effectue un calcul sur un ensemble de lignes
li\'ees \`a la ligne courante, \emph{sans r\'eduire le nombre de lignes}
du r\'esultat. Elle utilise la clause \texttt{OVER()}.
\end{definition}

\begin{theoreme}[Diff\'erence GROUP BY vs.\ Window Function]
\begin{itemize}
    \item \texttt{GROUP BY} + agr\'egation~: r\'eduit les lignes (une ligne par groupe).
    \item Fonction de fen\^etre~: conserve toutes les lignes tout en ajoutant
          une valeur calcul\'ee.
\end{itemize}
\end{theoreme}

\subsection{Syntaxe OVER}

\begin{formules}[title=\textbf{Syntaxe des fonctions de fen\^etre}]
\begin{minted}{sql}
fonction(args) OVER (
    [PARTITION BY colonnes]
    [ORDER BY colonnes [ASC|DESC]]
    [ROWS BETWEEN debut AND fin]
)
\end{minted}
\end{formules}

\subsection{Agr\'egation fen\^etr\'ee}

\begin{codeblock}[title=\textbf{Moyenne fen\^etr\'ee par fili\`ere}]
\begin{minted}{sql}
SELECT e.nom, e.filiere, i.note,
       AVG(i.note) OVER (PARTITION BY e.filiere) AS moy_filiere,
       i.note - AVG(i.note) OVER (PARTITION BY e.filiere) AS ecart
FROM Etudiants e
JOIN Inscriptions i ON e.id_etudiant = i.id_etudiant
ORDER BY e.filiere, e.nom;
\end{minted}
\end{codeblock}

\begin{codeoutput}
\begin{verbatim}
nom     | filiere | note | moy_filiere | ecart
--------+---------+------+-------------+------
Bernard | Info    | 15.5 | 13.1        | 2.4
Bernard | Info    | 13.0 | 13.1        | -0.1
Leroy   | Info    | 9.0  | 13.1        | -4.1
Petit   | Info    | 12.0 | 13.1        | -1.1
Petit   | Info    | 16.0 | 13.1        | 2.9
Moreau  | Maths   | 17.5 | 16.5        | 1.0
Moreau  | Maths   | 14.0 | 16.5        | -2.5
Roux    | Maths   | 18.0 | 16.5        | 1.5
\end{verbatim}
\end{codeoutput}

\subsection{ROW\_NUMBER, RANK, DENSE\_RANK}

\begin{definition}[Fonctions de classement]
\begin{itemize}
    \item \texttt{ROW\_NUMBER()} --- num\'ero s\'equentiel unique dans la partition.
    \item \texttt{RANK()} --- rang avec trous (ex-aequo re\c{c}oivent le m\^eme rang,
          le rang suivant est incr\'ement\'e du nombre d'ex-aequo).
    \item \texttt{DENSE\_RANK()} --- rang sans trous.
\end{itemize}
\end{definition}

\begin{codeblock}[title=\textbf{Classement des \'etudiants par note}]
\begin{minted}{sql}
SELECT e.nom, i.note,
       ROW_NUMBER() OVER (ORDER BY i.note DESC) AS row_num,
       RANK()       OVER (ORDER BY i.note DESC) AS rang,
       DENSE_RANK() OVER (ORDER BY i.note DESC) AS rang_dense
FROM Etudiants e
JOIN Inscriptions i ON e.id_etudiant = i.id_etudiant;
\end{minted}
\end{codeblock}

\begin{codeoutput}
\begin{verbatim}
nom     | note | row_num | rang | rang_dense
--------+------+---------+------+-----------
Roux    | 18.0 | 1       | 1    | 1
Moreau  | 17.5 | 2       | 2    | 2
Petit   | 16.0 | 3       | 3    | 3
Bernard | 15.5 | 4       | 4    | 4
Moreau  | 14.0 | 5       | 5    | 5
Bernard | 13.0 | 6       | 6    | 6
Petit   | 12.0 | 7       | 7    | 7
Leroy   | 9.0  | 8       | 8    | 8
\end{verbatim}
\end{codeoutput}

\begin{codeblock}[title=\textbf{Meilleur \'etudiant par fili\`ere (Top-N par groupe)}]
\begin{minted}{sql}
WITH Classement AS (
    SELECT e.nom, e.filiere, i.note,
           ROW_NUMBER() OVER (
               PARTITION BY e.filiere
               ORDER BY i.note DESC
           ) AS rn
    FROM Etudiants e
    JOIN Inscriptions i ON e.id_etudiant = i.id_etudiant
)
SELECT nom, filiere, note
FROM Classement
WHERE rn = 1;
\end{minted}
\end{codeblock}

\subsection{LAG, LEAD}

\begin{definition}[LAG et LEAD]
\begin{itemize}
    \item \texttt{LAG(col, n, defaut)} --- valeur de \texttt{col} $n$ lignes \emph{avant}
          la ligne courante.
    \item \texttt{LEAD(col, n, defaut)} --- valeur de \texttt{col} $n$ lignes \emph{apr\`es}
          la ligne courante.
\end{itemize}
\end{definition}

\begin{codeblock}[title=\textbf{Comparaison avec la ligne pr\'ec\'edente}]
\begin{minted}{sql}
SELECT e.nom, i.note,
       LAG(i.note, 1) OVER (ORDER BY i.note DESC) AS note_precedente,
       i.note - LAG(i.note, 1) OVER (ORDER BY i.note DESC) AS diff
FROM Etudiants e
JOIN Inscriptions i ON e.id_etudiant = i.id_etudiant
ORDER BY i.note DESC;
\end{minted}
\end{codeblock}

\subsection{Fen\^etres glissantes (Frame)}

\begin{codeblock}[title=\textbf{Moyenne mobile sur 3 lignes}]
\begin{minted}{sql}
SELECT e.nom, i.note,
       AVG(i.note) OVER (
           ORDER BY i.note
           ROWS BETWEEN 1 PRECEDING AND 1 FOLLOWING
       ) AS moyenne_mobile
FROM Etudiants e
JOIN Inscriptions i ON e.id_etudiant = i.id_etudiant
ORDER BY i.note;
\end{minted}
\end{codeblock}

\begin{formules}[title=\textbf{Sp\'ecifications de fen\^etre (Frame)}]
\begin{center}
\begin{tabular}{ll}
\toprule
\textbf{Sp\'ecification} & \textbf{Description} \\
\midrule
\texttt{ROWS BETWEEN UNBOUNDED PRECEDING AND CURRENT ROW} & Du d\'ebut \`a la ligne courante \\
\texttt{ROWS BETWEEN 1 PRECEDING AND 1 FOLLOWING} & Fen\^etre glissante de 3 lignes \\
\texttt{ROWS BETWEEN CURRENT ROW AND UNBOUNDED FOLLOWING} & De la ligne courante \`a la fin \\
\texttt{RANGE BETWEEN ...} & Bas\'e sur les valeurs (pas les positions) \\
\bottomrule
\end{tabular}
\end{center}
\end{formules}

\subsection{NTILE et calculs de pourcentage}

\begin{codeblock}[title=\textbf{Quartiles et pourcentage cumulatif}]
\begin{minted}{sql}
SELECT e.nom, i.note,
       NTILE(4) OVER (ORDER BY i.note DESC) AS quartile,
       ROUND(100.0 * SUM(i.note) OVER (ORDER BY i.note DESC
           ROWS BETWEEN UNBOUNDED PRECEDING AND CURRENT ROW)
           / SUM(i.note) OVER (), 1) AS pct_cumul
FROM Etudiants e
JOIN Inscriptions i ON e.id_etudiant = i.id_etudiant
ORDER BY i.note DESC;
\end{minted}
\end{codeblock}

\section{Int\'egration Python}

\begin{codeblock}[title=\textbf{Agr\'egation et window functions avec Python}]
\begin{minted}{python}
import sqlite3

conn = sqlite3.connect('universite.db')
cur = conn.cursor()

# Agregation avec GROUP BY
cur.execute('''
    SELECT e.nom, COUNT(*) as nb, ROUND(AVG(i.note), 2) as moy
    FROM Etudiants e
    JOIN Inscriptions i ON e.id_etudiant = i.id_etudiant
    GROUP BY e.id_etudiant, e.nom
    HAVING AVG(i.note) > 14
    ORDER BY moy DESC
''')
print("Moyennes > 14 :")
for row in cur.fetchall():
    print(f"  {row[0]:10} | {row[1]} cours | moyenne {row[2]}")

# Window function (SQLite >= 3.25)
cur.execute('''
    SELECT e.nom, i.note,
           RANK() OVER (ORDER BY i.note DESC) as rang
    FROM Etudiants e
    JOIN Inscriptions i ON e.id_etudiant = i.id_etudiant
''')
print("\nClassement :")
for nom, note, rang in cur.fetchall():
    print(f"  #{rang} {nom:10} : {note}")

conn.close()
\end{minted}
\end{codeblock}

\section*{Exercices}

\begin{exercice}
Pour chaque cours, affichez le nombre d'inscrits, la moyenne, la note
minimale et maximale. N'affichez que les cours ayant au moins 2 inscrits.
\end{exercice}

\begin{exercice}
Classez les \'etudiants par moyenne g\'en\'erale et affichez leur rang.
Utilisez \texttt{DENSE\_RANK} pour g\'erer les ex-aequo.
\end{exercice}

\begin{exercice}
Pour chaque inscription, affichez la note de l'\'etudiant, la moyenne
de sa fili\`ere, et l'\'ecart par rapport \`a cette moyenne.
Utilisez une fonction de fen\^etre.
\end{exercice}

\begin{exercice}
Affichez le top-2 des meilleures notes par cours. Utilisez
\texttt{ROW\_NUMBER} avec \texttt{PARTITION BY}.
\end{exercice}

\begin{exercice}
Calculez la somme cumulative des notes par \'etudiant, ordonn\'ee par
identifiant de cours, en utilisant une fen\^etre glissante.
\end{exercice}
